\subsection{Derivations in fields of characteristic zero}

THEOREM 1. --- Let K be a field of characteristic 0, L an extension of K and V a vector space over L. Let \(\Delta\) be a derivation of K into V, \((x_{i})_{i\in I}\) a transcendence basis of L over K and \((u_{i})_{i\in I}\) a family of elements of V. Then there exists a unique derivation D of L into V extending \(\Delta\) and such that \(\mathrm{D}(x_{i})=u_{i}\) for all \(i\in I\) .

Put \(E = K(x_{i})_{i \in I}\) ; Prop. 3 (V, p.~128) shows that \(\Delta\) extends to a unique derivation \(D_{0}\) of E into V such that \(D_{0}(x_{i}) = u_{i}\) for all \(i \in I\) . The field L is an algebraic extension of E and since L is of characteristic 0, L is separable over E (V, p.~37, Cor.). Therefore (V, p.~129, Prop. 4) \(D_{0}\) extends to a unique derivation D of L into V.

COROLLARY 1. --- Every derivation of K into V extends to a derivation of L into V.

COROLLARY 2. --- Let E be a subextension of L and U the vector sub-L-space of \(\Omega_{\mathrm{K}}(\mathrm{L})\) generated by the differentials of elements of E. For an element \(x\) to be algebraic over E it is necessary and sufficient that \(dx \in \mathrm{U}\) .

For each \(y \in L\) let \(\mathrm{D}(y)\) be the residue class of dy mod U. Then D is an E-derivation of L into \(\Omega_{\mathrm{K}}(\mathrm{L})/\mathrm{U}\) . Since K is of characteristic 0, every algebraic extension of E is separable (V, p.~37, Cor.) ; if \(x \in L\) is algebraic over E, we have Dx = 0 by Prop. 4 (V, p.~129), that is, \(dx \in U\) .

If \(x \in L\) is transcendental over E, there exists an E-derivation \(\Delta\) of \(\mathrm{E}(x)\) into L such that \(\Delta(x) = 1\) (V, p.~128, Prop. 3); by Th. 1, \(\Delta\) extends to an E-derivation D of L into L. Let \(\varphi\) be the linear form on \(\Omega_{\mathrm{K}}(\mathrm{L})\) such that \(D = \varphi \circ d\) ; we have \(\varphi(dy) = 0\) for \(y \in E\) and \(\varphi(x) = 1\) , whence \(dx \notin U\) .

COROLLARY 3. --- For an element x of L to be algebraic over K it is necessary and sufficient that dx = 0 in \(\Omega_{\mathrm{K}}(\mathrm{L})\) . In particular, for L to be an algebraic extension of K it is necessary and sufficient that \(\Omega_{\mathrm{K}}(\mathrm{L}) = 0\) .

This follows immediately from Cor. 2 on taking E = K.

COROLLARY 4. --- Let K, L and M be fields of characteristic 0 such that \(K \subset L \subset M\) ; then the canonical mapping \(\alpha: \Omega_{K}(L) \otimes_{L} M \to \Omega_{K}(M)\) is injective. This follows at once from Cor. 1 and V, p.~128, Prop. 2.

THEOREM 2.--- Let K be a field of characteristic 0, L an extension of K and \((x_{i})_{i\in I}\) a family of elements of L.

\begin{enumerate}
\def\labelenumi{\alph{enumi})}
\item
  For \((x_{i})_{i\in I}\) to be algebraically free over \(K\) it is necessary and sufficient that the differentials \(dx_{i}\) (for \(i\in I\) ) in \(\Omega_K(L)\) should be linearly free over \(L\) .
\item
  For L to be algebraic over \(\mathbf{K}(x_{i})_{i\in I}\) it is necessary and sufficient that the differentials \(dx_{i}\) for \(i\in I\) generate the vector space \(\Omega_{\mathbf{K}}(\mathbf{L})\) over L.
\item
  For \((x_{i})_{i\in I}\) to be a transcendence basis of L over K, it is necessary and sufficient that the family \((dx_{i})_{i\in I}\) should be a basis of \(\Omega_{\mathbf{K}}(\mathbf{L})\) over L.
\end{enumerate}

For every \(i \in I\) let \(E_i\) be the subfield \(\mathbf{K}(x_j)_{j \in I - \{i\}}\) of \(L\) . For the family \((x_i)_{i \in I}\) to be algebraically free over \(K\) it is necessary and sufficient (V, p.~108, Prop. 5) that \(x_i\) should be transcendental over \(E_i\) for each \(i \in I\) . By Cor. 2 of Th. 1 this means that for each \(i \in I\) the differential \(dx_i\) is not a linear combination with coefficients in \(L\) of the differentials \(dx_j\) with \(j \neq i\) in \(I\) ; this latter condition just means that the family \((dx_i)_{i \in I}\) is free over \(L\) , whence \(a\) ).

Assertion \(b)\) follows at once from Cor. 2 of Th. 1 and \(c)\) is a consequence of \(a)\) and \(b).\)

COROLLARY. --- We have \([\Omega_{\mathrm{K}}(\mathrm{L}):\mathrm{L}]=\mathrm{tr}\cdot\mathrm{deg}_{\mathrm{K}}\mathrm{L}\) when K is of characteristic 0.

